% ============================================================
% Capítulo 3 — Sistema descentralizado
% ============================================================
\chapter{Sistema descentralizado multicámara}
\label{cap:sistema-descentralizado}

Este capítulo describe el sistema completo de vigilancia en el que
se enmarca U-TextReIDNet. \textbf{No hay código en este proyecto
que lo implemente}, pero es útil entenderlo para saber qué parte
del paper es la que realmente vive en \texttt{textreid-train}.

\section{Arquitectura del sistema}

El sistema propuesto en \file{docs/paper.md} (§\textsc{Iii}-B)
tiene tres capas:

\begin{enumerate}
  \item \textbf{Nodos de cámara} ($d_1, \ldots, d_n$): capturan
        vídeo, mantienen galerías locales y ejecutan
        U-TextReIDNet para hacer inferencia local.
  \item \textbf{Infraestructura de comunicaciones}: red entre
        nodos y estación de control.
  \item \textbf{Estación de control}: servidor de aplicaciones
        (con U-TextReIDNet también) + aplicación web
        (\emph{Flask}).
\end{enumerate}

\begin{figure}[H]
\centering
\begin{tikzpicture}[
    font=\sffamily\small,
    node distance=0.5cm and 0.4cm,
    every node/.style={align=center},
    cam/.style={
        rectangle, draw=black, thick, rounded corners=2pt,
        minimum height=1.2cm, minimum width=1.8cm,
        fill=yellow!15
    },
    serv/.style={
        rectangle, draw=black, thick, rounded corners=2pt,
        minimum height=1.5cm, minimum width=3cm,
        fill=red!15
    },
    cli/.style={
        rectangle, draw=black, thick, rounded corners=2pt,
        minimum height=1.2cm, minimum width=2cm,
        fill=cyan!15
    },
    arr/.style={
        -{Stealth[length=2.5mm,width=2mm]}, thick
    }
]

% Nodos de cámara
\node[cam] (d1) {$d_1$\\\footnotesize Jetson};
\node[cam, right=0.5cm of d1] (d2) {$d_2$\\\footnotesize Jetson};
\node[cam, right=0.5cm of d2] (dn) {$d_n$\\\footnotesize Jetson};
\node[right=0.1cm of dn] {\ldots};

% Estación de control
\node[serv, below=2cm of d2] (srv) {\textbf{Servidor de apps}\\CPU i7 + GTX1050TI};

% Cliente
\node[cli, below=1.5cm of srv] (cli) {\textbf{Cliente}\\Navegador web};

% Conexiones
\draw[arr] (d1) -- (srv);
\draw[arr] (d2) -- (srv);
\draw[arr] (dn) -- (srv);
\draw[arr, dashed] (srv) -- (cli) node[midway, right] {\footnotesize Flask};

\end{tikzpicture}
\caption{Topología del sistema descentralizado. Cada cámara corre
U-TextReIDNet localmente; sólo transmite vídeo cuando su puntuación
supera un umbral y es la mejor de todas.}
\label{fig:topologia}
\end{figure}

\section{Flujo operativo}

El paper describe cinco fases (\file{docs/paper.md}, líneas 185--195):

\subsection*{a) Inicialización y solicitud}

Al encenderse, cada cámara se conecta al servidor de aplicaciones
y reporta su estado (\emph{online-idle}, \emph{online-búsqueda},
\emph{transmitiendo}, \emph{offline}). El usuario inicia o detiene
búsquedas desde la aplicación.

\subsection*{b) Captura de fotogramas y galería local}

Cuando una cámara recibe la orden de búsqueda, captura vídeo y
usa R-MHParsNet para generar una galería local de instancias
humanas detectadas en el fotograma actual. \emph{Cada fotograma
se trata como datos nuevos}: no hay relación temporal entre
galerías consecutivas.

\subsection*{c) Inferencia local}

Cada cámara calcula la similitud coseno entre cada imagen de su
galería local y la descripción textual proporcionada por el
usuario. Selecciona la imagen mejor clasificada y, si su
puntuación supera un umbral $\theta$, la procesa para incluir
un bounding box alrededor de la persona objetivo.

\subsection*{d) Identificación global}

Cada segundo, todas las cámaras cuyas puntuaciones superen
$\theta$ envían sus puntuaciones al servidor. El servidor
selecciona la puntuación más alta y notifica a esa cámara para
que transmita vídeo.

\subsection*{e) Finalización}

El usuario debe detener una búsqueda antes de iniciar otra.

\section{Implementación de la estación de control}

La aplicación de usuario es una webapp \emph{Flask}
\cite{grinberg2018flask} que corre en el servidor de aplicaciones.
Se divide en dos subsistemas:

\begin{itemize}
  \item \textbf{Subsistema de comunicación}: usa
        \href{https://zeromq.org/}{ZMQ} para mensajería
        asíncrona entre el servidor y los nodos, y
        \href{https://mqtt.org/}{MQTT} para publicar
        resultados de búsqueda desde los nodos al servidor.
  \item \textbf{Subsistema de vídeo}: usa
        \href{https://github.com/jeffbass/imagezmq}{ImageZMQ}
        para gestionar las transmisiones de vídeo entrantes
        cuando se detecta a la persona objetivo.
\end{itemize}

Estos subsistemas corren en hilos separados para que la UI no se
bloquee durante el procesamiento.

\section{Por qué este proyecto no incluye el sistema}

El repositorio \texttt{textreid-train} es una versión
\emph{training-only} del proyecto original. El README lo explica:

\begin{quote}
El proyecto original es un sistema completo de vigilancia
multicámara descentralizada con scripts de inferencia, una webapp
Flask, MySQL/MQTT y detección humana. \emph{Esta versión quita
todo eso} -- sólo tiene lo necesario para \textbf{entrenar} y
\textbf{evaluar}.
\end{quote}

\noindent Las piezas eliminadas son:

\begin{itemize}
  \item \file{model/unity.py}: wrapper de detección + re-ID
        (innecesario para entrenamiento puro).
  \item \file{model/human_detection_network.py}: detector
        (R-MHParsNet no se entrena aquí).
  \item \file{model/feature_pyramid_network.py}: usado solo por
        el detector.
  \item \file{model/mask_head.py}, \file{parsing_head.py}:
        cabezas del detector.
  \item \file{surveillance_application/}: la webapp Flask.
  \item \file{demo/}: demo web.
  \item \file{inference/}: scripts de inferencia sobre
        imágenes/vídeos.
\end{itemize}

Para desplegar el sistema completo habría que integrar este modelo
con un detector (YOLOv8, por ejemplo) y la infraestructura
descrita arriba. Eso queda fuera del alcance del libro.

\section{Rendimiento reportado en el paper}

El paper mide tiempos de inferencia en dos plataformas:

\begin{table}[H]
\centering
\caption{Tiempos de inferencia reportados en el paper (media sobre
20 iteraciones).}
\label{tab:tiempos-paper}
\begin{tabular}{lrr}
\toprule
\textbf{Modelo} & \textbf{PC (CPU+GPU)} & \textbf{Jetson Nano} \\
\midrule
MFPE        & 0{,}69204 s & --- \\
TextReIDNet & 0{,}00477 s & --- \\
U-TextReIDNet & 0{,}04995 s & 0{,}75013 s \\
\bottomrule
\end{tabular}
\end{table}

\noindent Para una galería de 10 personas:

\begin{table}[H]
\centering
\caption{Tiempos de inferencia para 10 imágenes de galería.}
\begin{tabular}{lrr}
\toprule
\textbf{Modelo} & \textbf{PC} & \textbf{Jetson Nano} \\
\midrule
MFPE          & 2{,}46102 s & --- \\
TextReIDNet   & 1{,}28512 s & --- \\
U-TextReIDNet & 0{,}09525 s & 0{,}75013 s \\
\bottomrule
\end{tabular}
\end{table}

\noindent La conclusión clave es que U-TextReIDNet puede procesar
galerías pequeñas en tiempo real incluso en hardware embebido, y
su tiempo de inferencia escala linealmente con el número de
personas ($O(N \cdot n)$ donde $N$ es la densidad y $n$ el
número de características).

\section{Lo que este proyecto sí hace}

Aunque no incluye el sistema distribuido, el modelo entrenado en
\texttt{textreid-train} \emph{es} el mismo TextReIDNet que se usa
en el paper como submódulo de U-TextReIDNet. Por tanto, un
checkpoint producido aquí puede integrarse en el sistema
completo con sólo cambiar la arquitectura de detección.
